%% ece448_labpacket.tex -- shared macros for every ECE 448 lab packet.
%%
%% A packet \input's this after exam_config and before \begin{document}, with
%% \ifkey already set (\newif\ifkey, then \keytrue or \keyfalse). The L5 Lab,
%% ECE448_Lab_L05_FirstFlowgraphs.tex, is the worked example; the other labs
%% load this file as each is reworked (ROADMAP, "Labs follow 444's lab rules").
%%
%% The shape (VOICE.md 34, COURSE_SPEC "Lab packet structure" in ECE 444):
%%   * Predict, Measure, Analyze, one submission, about four pages.
%%   * Predicted sits beside measured: each table appears once, with a
%%     Predicted and a Measured column per quantity, or a Predicted row above
%%     the Measured row where pairs do not fit.
%%   * A predicted cell is shaded with a gray box drawn INSIDE the cell, never
%%     \cellcolor: \cellcolor covers the table rules, so a run of predicted
%%     cells reads as one gray block on screen. It is named p1.<the Part 2
%%     field name>. A cell with no prediction shows a dash (\nopred).
%%   * Answer boxes are \boxw (1.5 in) wide, right-justified in one 0.42
%%     column (anscol) beside a 0.55 worked column, or in one row along the
%%     bottom of the part.
%%
%% Fields. In the blank copy every predicted and measured cell is also an
%% AcroForm text field of the cell's name, so a cadet can type into the PDF,
%% and the lab page's entry boxes (records.js) write into the Measured fields
%% by the same names (scripts/verify/check_records.py checks the names). The
%% KEY prints the expected value in red instead. Wrap the fields in ONE
%% \begin{Form} ... \end{Form} (hyperref) around the packet body.
%%
%% Answers. No answer is in a public file. Every expected value is
%% \keysol{<id>}, defined in the private ECE448-faculty repo at
%% lab-keys/src/<packet>_solutions.tex with \defkeysol{<id>}{<text>}, and
%% loaded by \loadkeysols{<packet>} (a KEY build without the file stops).

\RequirePackage{array}
\RequirePackage{xcolor}

% ---- answers from the private repo --------------------------------------
\long\def\defkeysol#1#2{\expandafter\long\expandafter\def\csname keysol@#1\endcsname{#2}}
\def\keysol#1{\ifcsname keysol@#1\endcsname\csname keysol@#1\endcsname\fi}
\newcommand{\loadkeysols}[1]{%
	\ifkey
		\IfFileExists{#1_solutions.tex}%
			{\input{#1_solutions}}%
			{\errmessage{KEY build needs #1_solutions.tex from ECE448-faculty/lab-keys/src (set KEYDIR)}}%
	\fi}

% ---- sizes and colors ----------------------------------------------------
\newlength{\boxw}      \setlength{\boxw}{1.5in}    % every answer box
\newlength{\cellh}     \setlength{\cellh}{0.30in}  % a table cell's writable height
\newlength{\predinset} \setlength{\predinset}{2.5pt} % gap between the gray box and the rules
\definecolor{predgray}{gray}{0.86}
\definecolor{keyred}{rgb}{0.80,0.0,0.0}

% Fixed-width centered column for fillable cells.
\newcolumntype{K}[1]{>{\centering\arraybackslash}p{#1}}

% The 0.55 worked column and the 0.42 answer column (VOICE.md 34).
\newenvironment{workcol}{\begin{minipage}[t]{0.55\linewidth}}{\end{minipage}}
\newenvironment{anscol}{\begin{minipage}[t]{0.42\linewidth}\raggedleft}{\end{minipage}}

% ---- the field itself ----------------------------------------------------
% \lp@field{<name>}{<width>}{<height>}: a borderless, transparent text field
% in the blank copy; nothing in the KEY.
\makeatletter
\newcommand{\lp@field}[3]{%
	\ifkey\else
		\TextField[name={#1}, width=#2, height=#3, borderwidth=0,
			bordercolor={}, backgroundcolor={}, charsize=9pt, align=1]{}%
	\fi}
% \lp@cell{<width>}{<height>}{<name>}{<key text>}{<shade:0|1>}
\newcommand{\lp@cell}[5]{%
	\begingroup\setlength{\fboxsep}{0pt}%
	\ifnum#5=1\colorbox{predgray}{\fi
	\begin{minipage}[c][#2][c]{#1}\centering
		\ifkey{\color{keyred}\small #4}\else\lp@field{#3}{#1}{#2}\fi
	\end{minipage}%
	\ifnum#5=1}\fi
	\endgroup}

% ---- predicted cells (shaded) ---------------------------------------------
% \predv[<Part 2 field>]{<key>}  in a table cell: the gray box fills the
%                                 cell's width less \predinset on each side.
% \predh[<Part 2 field>]{<key>}  a \boxw box for an answer column or a row of
%                                 boxes along the bottom of a part.
% \predl[<Part 2 field>]{<width>}{<key>}  a long predicted box of any width,
%                                 for an entry that does not fit \boxw.
% Each is named p1.<field>, so the Predicted cell of a.t20 is p1.a.t20.
\newcommand{\predv}[2][]{%
	\rule[-0.5\cellh]{0pt}{\dimexpr\cellh+2\predinset\relax}%
	\lp@cell{\dimexpr\linewidth-2\predinset\relax}{\cellh}{p1.#1}{#2}{1}}
\newcommand{\predh}[2][]{\lp@cell{\boxw}{\cellh}{p1.#1}{#2}{1}}
\newcommand{\predl}[3][]{\lp@cell{#2}{\cellh}{p1.#1}{#3}{1}}

% ---- measured cells (white) ------------------------------------------------
% \kv[<field>]{<key>}   in a table cell; \kvh[<field>]{<key>}, a \boxw box.
\newcommand{\kv}[2][]{%
	\rule[-0.5\cellh]{0pt}{\dimexpr\cellh+2\predinset\relax}%
	\lp@cell{\dimexpr\linewidth-2\predinset\relax}{\cellh}{#1}{#2}{0}}
\newcommand{\kvh}[2][]{\fbox{\lp@cell{\boxw}{\cellh}{#1}{#2}{0}}}

% A cell with no prediction.
\newcommand{\nopred}{--}

% ---- a screenshot box -------------------------------------------------------
% \shotbox{<height>}{<caption>}{<key>}: a framed region the full text width to
% paste a screenshot into, captioned below. The KEY prints what it should show.
\newcommand{\shotbox}[3]{%
	\par\noindent\fbox{\begin{minipage}[c][#1][c]{\dimexpr\linewidth-2\fboxsep-2\fboxrule\relax}\centering
		\ifkey{\color{keyred}\small #3}\else{\color{gray}\small paste here}\fi
	\end{minipage}}\par\nopagebreak[4]
	\vspace{2pt}{\small #2}\par}

% \shotpane{<width>}{<height>}{<label>}{<key>}: one labeled pane of a figure
% made of several screenshots side by side (e.g. three settings of one
% control); set the panes in a row with \hfill, then the caption.
\newcommand{\shotpane}[4]{%
	\begin{minipage}[t]{#1}\centering
		\fbox{\begin{minipage}[c][#2][c]{\dimexpr\linewidth-2\fboxsep-2\fboxrule\relax}\centering
			\ifkey{\color{keyred}\footnotesize #4}\else{\color{gray}\small paste here}\fi
		\end{minipage}}\\[2pt]{\small #3}
	\end{minipage}}

% \freqstrip{<half-span>}{<label step>}{<caption>}: a blank frequency axis
% from -<half-span> to +<half-span> to mark predictions on (kHz). It prints
% no answer in either copy.
\newcommand{\freqstrip}[3]{%
	\begin{center}\setlength{\unitlength}{\dimexpr0.8\linewidth/(2*#1)\relax}%
	\begin{picture}(\numexpr2*#1\relax,6)(-#1,-4)
		\put(-#1,0){\line(1,0){\numexpr2*#1\relax}}
		\multiput(-#1,0)(#2,0){\numexpr2*#1/#2+1\relax}{\line(0,1){1.2}}
		\count255=-#1\relax
		\loop\ifnum\count255>#1\else
			\put(\the\count255,-2.6){\makebox(0,0){\footnotesize\ifnum\count255>0 +\fi\the\count255}}%
			\advance\count255 by #2\relax
		\repeat
	\end{picture}\\[-2pt]{\small #3}\end{center}}

%% The procedure at a glance (Neil, 2026-11-22: "That way you can read the lab
%% packet and still have an idea of what the flow is"). A framed box on page 1
%% that lists the lab page's procedure blocks, one line each, with their step
%% ranges. It is the flow, not the steps: the lab page holds the full procedure.
%%   \begin{procbox}
%%     \procstep{1--7}{Build the tone flowgraph ...}
%%   \end{procbox}
\newenvironment{procbox}{%
	\par\smallskip\noindent\begin{lrbox}{0}\begin{minipage}{\dimexpr\linewidth-2\fboxsep-2\fboxrule\relax}%
	\textbf{The procedure at a glance} \hfill {\small (full steps on the lab page)}\par\smallskip
	\begin{tabular}{@{}p{0.13\linewidth}p{0.83\linewidth}@{}}}%
	{\end{tabular}\end{minipage}\end{lrbox}\fbox{\usebox{0}}\par\smallskip}
\newcommand{\procstep}[2]{\textbf{Steps #1} & #2 \\[1pt]}

%% The documentation statement (Neil, 2026-11-22): the last thing in every
%% packet, at the bottom of its last page, per the Dean's documentation policy
%% (course-policies.tex in latex-tools). A fillable field in the blank copy.
%% Put \docstatement after the last question, before \end{Form}.
\newcommand{\docstatement}{%
	\par\vfill\noindent
	\fbox{\begin{minipage}{\dimexpr\linewidth-2\fboxsep-2\fboxrule\relax}
	\textbf{Documentation:} {\small List every source of help on this packet: people,
	websites, generative AI, with what each helped with. If you had none, write
	``Documentation: None.'' A packet without this statement is incomplete.}\par
	\smallskip
	\ifkey{\color{keyred}\small (The cadet's own statement.)}\par\vspace{0.5in}%
	\else\TextField[name={doc}, multiline=true, width=\linewidth, height=0.65in,
		borderwidth=0, bordercolor={}, backgroundcolor={}, charsize=9pt]{}\fi
	\end{minipage}}\par}
\makeatother
